% paper.tex — PyPSA-Libya dissemination paper skeleton
% Structural template follows Gallego-Castillo & Victoria (2025),
% "PyPSA-Spain: An extension of PyPSA-Eur to model the Spanish energy system",
% Energy Strategy Reviews 60, 101764.
%
% Requires elsarticle.cls (download from https://ctan.org/pkg/elsarticle
% and place it in this folder — it is an Elsevier-distributed class file
% and is not reproduced here).

\documentclass[preprint,12pt]{elsarticle}

\usepackage{graphicx}
\usepackage{amsmath}
\usepackage{hyperref}

\journal{Energy Strategy Reviews}  % or your intended target journal

\begin{document}

\begin{frontmatter}

\title{PyPSA-Libya: An extension of PyPSA-Earth to model the Libyan energy system}

\author[a]{Mohamed Abuella}
\affiliation[a]{ Email= {mhdabuella@gmail.com}}

% \affiliation[a]{organization={Independent Researcher},
%                  city={Newcastle upon Tyne},
%                  country={UK}}

\begin{abstract}
This work presents PyPSA-Libya, an open-source model of the Libyan electricity
system, developed as a national-scale fork of PyPSA-Earth. [Fill in: spatial
resolution, key data sources used/available, headline case-study finding.]
\end{abstract}

\begin{keyword}
Energy system modelling \sep Libya \sep Open-source \sep PyPSA
\end{keyword}

\end{frontmatter}

\section{Introduction}
% Mirror PyPSA-Spain §1 structure: (i) why decarbonisation/grid modelling
% matters for this country, (ii) what PyPSA-Earth already provides,
% (iii) limitations of the continental model for national planning,
% (iv) what a national fork adds, (v) paper roadmap.


Test bibliography ~\cite{almaktar2021prospects, eljrushi1990solar, fakhreddine2025global, terlouw2022large, fattahi2024opportunities, abdelkhalek2025pypsa, gallego2025pypsaspain, johnson2025realistic, alanazi2025renewable, rahil2018techno, achhammer2026impact}


\section{Novel functionalities developed in PyPSA-Libya}

\subsection{Interconnection modelling}
% Tunisia and Egypt border-node + virtual-generator pattern
% (PyPSA-Spain §2.3 equivalent). Optional Italy HVDC scenario switch.

\subsection{Technology cost assumptions}
% data_LY/costs/costs_updates.csv — PyPSA-Spain §2.4.3 equivalent.

\subsection{Data limitations}
% Be explicit: no Datadis/ESIOS equivalent for Libya. State what proxy
% data was used instead (e.g. global GDP/population rasters, any GECOL
% figures obtained) and the resulting uncertainty.

\section{Case study}
% Mirror PyPSA-Spain §3–4: define scenario, target year, node-count
% sensitivity sweep, compare optimal mix against any available national
% plan/target (if one exists for Libya) or against current installed
% capacity as a sanity baseline.

\begin{figure}[ht]
\centering
    \includegraphics[width=0.85\textwidth]{Figures/Libya_H2/Energy Carriers Analysis/Scenario_A_physical_generation_by_carrier_filtered.png}
\caption{Capacity Scenario A}
    \label{fig:sect_coupl_capacity_sen_a}
\end{figure}

\section{Conclusions}

\section*{Code availability}
The code of PyPSA-Libya is available at
\url{https://github.com/<you>/pypsa-libya}.

\section*{Data availability}
Pre-computed inputs and results are archived at [Zenodo/OSF DOI].


%=========================================
% \begin{thebibliography}{00}
\bibliographystyle{elsarticle-num}
\bibliography{Libya_H2}
% \end{thebibliography}

\end{document}


\maketitle

\section{Introduction}

\end{document}
